\documentclass{article}
\usepackage{amsmath}
\usepackage{graphicx}
\usepackage{booktabs}
\usepackage{pgfplots}
\pgfplotsset{compat=1.16}

\title{The Impact of AI Augmentation on Quality and Defect Severity: An Economic Analysis}
\author{David Adams}
\date{\today}

\begin{document}

\maketitle

\section{Introduction}
This study evaluates the impact of AI augmentation on the quality of work, focusing on the reduction of defects and the economic implications of defect severity. We compare scenarios with and without AI assistance, quantifying improvements in quality and cost savings.

\section{Methodology}
We model the production of 1,000 units under two conditions: with AI augmentation and without AI (manual work). Defects are categorized by severity (low, medium, high), and costs are assigned based on the severity of each defect type.

\subsection{Defect Severity and Costs}
Defects are classified into three severity levels:
\begin{itemize}
    \item \textbf{Low severity}: Cosmetic or minor issues (cost: \$10 per defect).
    \item \textbf{Medium severity}: Functional but non-critical issues (cost: \$50 per defect).
    \item \textbf{High severity}: Critical or safety-related issues (cost: \$200 per defect).
\end{itemize}

\subsection{Data}
The number of defects for each severity level, with and without AI, is shown in Table \ref{table:defects}.

\begin{table}[h]
\centering
\caption{Defect Counts by Severity}
\label{table:defects}
\begin{tabular}{lcc}
\toprule
Severity      & With AI & Without AI \\
\midrule
Low           & 20      & 50         \\
Medium        & 20      & 70         \\
High          & 10      & 30         \\
\bottomrule
\end{tabular}
\end{table}

\section{Results}
\subsection{Defect Costs}
The total cost of defects is calculated as:
\[
\text{Total Defect Cost} = \sum (\text{Number of Defects} \times \text{Cost per Defect})
\]

With AI:
\[
\text{Total Defect Cost} = (20 \times 10) + (20 \times 50) + (10 \times 200) = 200 + 1000 + 2000 = \$3,200
\]

Without AI:
\[
\text{Total Defect Cost} = (50 \times 10) + (70 \times 50) + (30 \times 200) = 500 + 3500 + 6000 = \$10,000
\]

\subsection{Net Savings}
The total cost with AI includes an additional \$1,000 for AI implementation:
\[
\text{Total Cost with AI} = \$3,200 + \$1,000 = \$4,200
\]

The net savings with AI is:
\[
\text{Net Savings} = \$10,000 - \$4,200 = \$5,800 \text{ per 1,000 units}
\]

\subsection{Defect Reduction}
AI reduces the number of defects across all severity levels:
\begin{itemize}
    \item Low severity: 60\% reduction (from 50 to 20).
    \item Medium severity: 71.4\% reduction (from 70 to 20).
    \item High severity: 66.7\% reduction (from 30 to 10).
\end{itemize}

\begin{figure}[h]
\centering
\begin{tikzpicture}
\begin{axis}[
    ybar,
    ylabel=Number of Defects,
    symbolic x coords={Low, Medium, High},
    xtick=data,
    nodes near coords,
    nodes near coords align={vertical},
    title=Defect Severity: AI vs. Manual,
    legend pos=north west,
]
\addplot coordinates {(Low,20) (Medium,20) (High,10)};
\addplot coordinates {(Low,50) (Medium,70) (High,30)};
\legend{With AI, Without AI}
\end{tikzpicture}
\end{figure}

\section{Economic Implications}
\subsection{Cost Savings}
AI augmentation results in significant cost savings by reducing defects, particularly high-severity defects. The net savings of \$5,800 per 1,000 units justifies the \$1,000 cost of AI implementation.

\subsection{Quality Improvement}
The reduction in defects improves overall product quality, leading to higher customer satisfaction, fewer recalls, and lower rework costs.

\section{Conclusion}
AI augmentation significantly improves quality by reducing defects of all severity levels. The economic benefits, particularly from reducing high-severity defects, make AI a valuable investment for industries where quality is critical.

\section{Discussion}
Future work could explore:
\begin{itemize}
    \item Tuning AI to further reduce high-severity defects.
    \item Long-term learning effects of AI on defect reduction.
    \item Combining AI with worker training for enhanced quality improvements.
\end{itemize}

\end{document}
